%%%%%解析几何模版%%%%%
\chapter*{解析几何每日一题}

\refer{《高等数学》同济版}

\exercise{求通过$x$轴和点$\qty(4,-3,-1)$的平面方程}

\teacher{林经翔}

\postdate{2023年4月18日}

\result{y-3z=0}

\solution{
    设平面的一般方程为
    \formula{Ax+By+Cz+D=0}
    由于平面通过$x$轴，故它的法向量垂直于$x$轴，即$A=0$.

    又因为平面通过$x$轴，它必通过原点，于是$D=0$.

    因此平面的方程简化为
    \formula{By+Cz=0}
    因为平面通过点$\qty(4,-3,1)$，代入方程得
    \formula{-3B-C=0}
    取$B=1,\ C=-3$，则所求平面方程为
    \formula{y-3z=0}
}

\point{
    平面的一般方程
    
    1.平面的一般方程为$Ax+by+Cz=0$，法线向量$\vv{n}=\qty(A,B,C)$

    2.熟悉一些特殊的方程和它们对应的图形特点
    \begin{enumerate}
        \item 当$D=0$时，$Ax+By+Cz$表示一个平行于（或包含）$x$轴得平面
        \item 当$A=0$时，法向量$\vv{n}=\qty(0,B,C)$垂直于$x$轴.
        
        \hspace{2.2em} 此时，$Cz+D=0$表示一个平行于（或重合于）$xOy$面的平面

        \hspace{2.2em} 同样，$Ax+D=0$表示一个平行于（或重合于）$yOz$面的平面

        \hspace{2.2em} \qquad\quad $By+D=0$表示一个平行于（或重合于）$xOz$面的平面
    \end{enumerate}
}